\section{All rational parameters on the target line}\label{sec:local-global}

The full fiber is controlled by the roots of
\begin{equation}
 f_a(X)=(X^3-a)(X^2+X+1).
 \label{eq:rational-polynomial}
\end{equation}
For \(a\ne0,1\), its two factors are separable and coprime, so the
fiber is the root scheme described in \eqref{eq:main-fiber}.
The two factors cover different completions: the quadratic supplies a
root at primes \(1\pmod3\), and the cubic must supply one at every
remaining prime.  This gives a converse to the sufficient conditions
defining \(\mathcal A\).

\begin{theorem}[Rational classification]\label{thm:rational-classification}
For \(a\in\Q\), the fiber \(\Phi^{-1}(y_a)\) has points over every
completion of \(\Q\) and has no rational point if and only if
\(a\in\Q^\times\setminus\Q^{\times3}\) and
\begin{equation}\label{eq:rational-local}
 \begin{gathered}
 v_p(a)\equiv0\pmod3
       \quad\hbox{for every prime }p\equiv2\pmod3,\\
 a\in\Q_3^{\times3}.
 \end{gathered}
\end{equation}
Writing \(a=3^e u\) with \(u\in\Z_3^\times\), the second condition is
equivalent to
\begin{equation}
 e\equiv0\pmod3,\qquad u\equiv1\ \hbox{or }-1\pmod9.
 \label{eq:three-adic-cube}
\end{equation}
Equivalently, the Hasse-failing parameters are exactly those admitting a
representation
\begin{equation}
 a=\varepsilon b^3\frac{r}{s},
 \qquad \varepsilon\in\{1,-1\},\quad b\in\Q_{>0},
 \label{eq:rational-normal-form}
\end{equation}
where \(r,s\) are coprime positive cube-free integers, every prime
dividing \(rs\) is \(1\pmod3\), \(r\equiv s\pmod9\), and \(rs>1\).
\end{theorem}

\begin{proof}
First suppose \(a\ne0,1\).  The quadratic has no rational root because
its discriminant is \(-3\).  The cubic has a rational root exactly when
\(a\) is a rational cube.  Thus the full fiber has no rational point
exactly when \(a\) is not a cube.

At a prime \(p\ne3\), the quadratic has a root in \(\Q_p\) exactly when
\(\Q_p\) contains a primitive cube root of unity, equivalently when
\(p\equiv1\pmod3\).  To see the converse as well as the forward
implication, its roots are integral units; their reductions are
nontrivial cube roots of unity, since \(X^2+X+1\) takes the value \(3\)
at \(1\).  Conversely, a nontrivial cube root of unity in
\(\F_p^\times\) is a simple quadratic root and lifts by Hensel's lemma.
This includes \(p=2\) in the no-root case.

If \(p\equiv2\pmod3\), cubing is an automorphism of
\(\F_p^\times\).  Every \(p\)-adic unit therefore has a cube root, by
Hensel's lemma applied to \(X^3-u\); its derivative at a nonzero
residue root is a unit.  It follows that \(a\) is a cube in \(\Q_p\)
if and only if \(3\mid v_p(a)\).  This proves the necessity and
sufficiency of all the inert-prime conditions.

At \(3\), the quadratic has no root because its discriminant \(-3\)
has odd valuation.  Hence the cubic must have a root, which is precisely
\(a\in\Q_3^{\times3}\).  A cube has valuation divisible by three, and
the cube of any \(3\)-adic unit is \(1\) or \(-1\pmod9\).
These necessary conditions are sufficient.  For a unit \(u=1+9k\),
the substitution \(X=1+3T\) gives
\[
 X^3-u=9\bigl(T+3T^2+3T^3-k\bigr).
\]
The expression in parentheses has derivative \(1\pmod3\) and has
the root \(T\equiv k\pmod3\), so Hensel's lemma gives a cube root
of \(u\).  If \(u\equiv-1\pmod9\), apply this argument to \(-u\)
and change the sign of the root.  The real completion always has the
cubic root \(\sqrt[3]{a}\).  This proves \eqref{eq:rational-local}.

For the normal form, write the reduced absolute fraction for \(a\) and
remove the largest cube from its numerator and denominator.  The
valuation conditions leave cube-free parts \(r,s\) supported only on
primes \(1\pmod3\); all remaining factors, including the sign, are
rational cubes.  Every such prime lies in
\[
 H=\{1,4,7\}\subset(\Z/9\Z)^\times.
\]
The cube of a \(3\)-adic unit is \(1\) or \(-1\pmod9\), whereas
\(H\cap\{1,-1\}=\{1\}\).  The unit condition at \(3\) is therefore
equivalent to \(r/s\equiv1\pmod9\), or \(r\equiv s\pmod9\).
Finally, coprime cube-free \(r,s\) give a rational cube exactly when
\(r=s=1\).  Conversely, a representation with the stated conditions
satisfies \eqref{eq:rational-local} and is not a rational cube.

It remains to check the excluded values, since they also lie on the
target line.  The inverse-coordinate calculation in
Proposition~\ref{prop:gauge} applies to every root for which \(E'\) is
a unit, and every source point satisfies
\(D=(16/9)E'=t^{-1}\).  Thus, without a squarefreeness assumption, the
fiber over a target with \(\pi\ne0\) is
\[
 \Spec\Q[S,E'(S)^{-1}]/(E(S)).
\]
At \(a=0\), the factor \(X^3\) is removed by this localization and
the remaining fiber is \(\Spec\Q[X]/(X^2+X+1)\), which has no real
point.  At \(a=1\),
\[
 f_1(X)=(X-1)(X^2+X+1)^2,
\]
and only the simple root \(X=1\) remains, giving a rational point.
Neither value is a Hasse failure.
\end{proof}

\begin{corollary}\label{cor:sufficient-family}
Every \(a\in\mathcal A\) gives an everywhere locally soluble fiber
without a rational point.
\end{corollary}

\begin{proof}
The prime-support condition gives \(v_p(a)=0\) at every prime
\(p\equiv2\pmod3\), and \(a\equiv1\pmod9\) makes \(a\) a
\(3\)-adic cube.  The noncube condition in the definition of
\(\mathcal A\) completes the application of
Theorem~\ref{thm:rational-classification}.
\end{proof}

For example, \(a=152=2^3\cdot19\) satisfies the criterion even though
its prime support includes \(2\), and \(a=19/37\) satisfies it since
\(19\equiv37\equiv1\pmod9\).  The earlier integer examples
\(a=19\) and \(a=91=7\cdot13\) also satisfy the full criterion.

\begin{remark}
The reducibility of \(f_a\) is forced by the local condition.  A connected
finite \'etale \(\Q\)-scheme of degree greater than one has a transitive
Galois action containing a derangement.  Chebotarev then gives an unramified
prime at which the scheme has no point.  Hasse-failing zero-dimensional
schemes must consequently have at least two arithmetic components; see
Sonn \cite{Sonn} for the corresponding subgroup-covering formulation.
\end{remark}

\begin{remark}
The rational classification above is a written argument.  The
machine-checked arithmetic certificate described in
Section~\ref{sec:verification} retains its stated scope, namely the
sufficient integer family \(\mathcal A\).
\end{remark}
